\documentclass{article}
\usepackage{amsmath,amssymb,amsthm}

\newcommand{\R}{\mathbb{R}}
\newcommand{\ip}[2]{\langle #1, #2 \rangle}
\newcommand{\half}{\tfrac{1}{2}}
\newcommand{\diag}{\operatorname{diag}}

\title{Geodesic Newton Direction for the Extended Embedding}
\author{}
\date{}

\begin{document}
\maketitle

\section{Setup}

Consider the extended embedding of a conic program with cone $\mathcal{C} = \R^n_+$:
\begin{align}
  \min \quad & \alpha \theta \notag \\
  \text{s.t.} \quad & Ax - b\tau = r_p \theta \label{eq:primal} \\
  & -A^T y - s + c\tau = r_d \theta \label{eq:dual} \\
  & \ip{b}{y} - \ip{c}{x} - \kappa = r_g \theta \label{eq:gap} \\
  & \ip{r_p}{y} + \ip{r_d}{x} + r_g \tau = -\alpha \label{eq:norm} \\
  & x \in \R^n_+, \; s \in \R^n_+, \; \tau \geq 0, \; \kappa \geq 0 \notag
\end{align}
where $A \in \R^{m \times n}$, $b \in \R^m$, $c \in \R^n$, and $y \in \R^m$, $\theta \in \R$ are free variables.

\section{Self-duality}

The embedding is self-dual: if $(x, s, y, \tau, \kappa, \theta)$ is feasible, then swapping
\[
  (x, s, \tau, \kappa, y) \;\longleftrightarrow\; (s, x, \kappa, \tau, -y)
\]
gives a dual-feasible point. This symmetry is the key structural property.

\section{Geodesic parametrization}

For the product cone $\R^n_+ \times \R^n_+ \times \R_+ \times \R_+$ with rank $R = 2n + 2$, the geodesic parametrization at scaling $W = \diag(w_x, w_s, w_\tau, w_\kappa)$ and barrier parameter $k = 1/\sqrt{\mu}$ gives:
\begin{align}
  \text{slack}_i &= \frac{1 - d_i}{k \, W_i}, \label{eq:slack} \\
  \lambda_i &= k \, W_i \, (1 + d_i). \label{eq:dual_var}
\end{align}

For the $x \geq 0$ constraint, the ``slack'' is $x$ itself and $\lambda_x$ is its dual. For $s \geq 0$, the slack is $s$ and $\lambda_s$ is its dual.

\section{Self-duality implies $d_x = -d_s$}

By self-duality, $x$ is simultaneously:
\begin{enumerate}
  \item The slack of the $x \geq 0$ constraint: $x_i = \frac{1 - d_{x,i}}{k \, w_{x,i}}$.
  \item The dual variable of the $s \geq 0$ constraint: $x_i = k \, w_{s,i} \, (1 + d_{s,i})$.
\end{enumerate}
Symmetrically, $s$ is:
\begin{enumerate}
  \item The slack of $s \geq 0$: $s_i = \frac{1 - d_{s,i}}{k \, w_{s,i}}$.
  \item The dual of $x \geq 0$: $s_i = k \, w_{x,i} \, (1 + d_{x,i})$.
\end{enumerate}

\subsection{Nesterov--Todd scaling}

Impose the NT scaling $w_{x,i} = w_{s,i}^{-1}$. Then from identifications (1) and (2) for $x$:
\[
  \frac{1 - d_{x,i}}{k \, w_{x,i}} = k \, w_{s,i} \, (1 + d_{s,i}) = k \, w_{x,i}^{-1} \, (1 + d_{s,i}).
\]
Multiply both sides by $k \, w_{x,i}$:
\[
  1 - d_{x,i} = k^2 \, (1 + d_{s,i}).
\]
Similarly from $s$:
\[
  \frac{1 - d_{s,i}}{k \, w_{s,i}} = k \, w_{x,i} \, (1 + d_{x,i}).
\]
Multiply by $k \, w_{s,i} = k \, w_{x,i}^{-1}$:
\[
  1 - d_{s,i} = k^2 \, (1 + d_{x,i}).
\]
Adding both equations:
\[
  (1 - d_{x,i}) + (1 - d_{s,i}) = k^2 \bigl[(1 + d_{s,i}) + (1 + d_{x,i})\bigr].
\]
\[
  2 - (d_{x,i} + d_{s,i}) = k^2 \bigl[2 + (d_{x,i} + d_{s,i})\bigr].
\]
Let $\sigma = d_{x,i} + d_{s,i}$:
\[
  2 - \sigma = k^2(2 + \sigma) \implies \sigma(1 + k^2) = 2(1 - k^2) \implies \sigma = \frac{2(1-k^2)}{1+k^2}.
\]
From the first equation alone:
\[
  1 - d_{x,i} = k^2(1 + d_{s,i}) \implies d_{x,i} = 1 - k^2 - k^2 d_{s,i}.
\]
From the second:
\[
  d_{s,i} = 1 - k^2 - k^2 d_{x,i}.
\]
Substituting:
\[
  d_{s,i} = 1 - k^2 - k^2(1 - k^2 - k^2 d_{s,i}) = (1-k^2)(1 - k^2) + k^4 d_{s,i}.
\]
\[
  d_{s,i}(1 - k^4) = (1-k^2)^2 \implies d_{s,i} = \frac{(1-k^2)^2}{(1-k^2)(1+k^2)} = \frac{1 - k^2}{1 + k^2}.
\]
This gives $d_{s,i}$ as a \emph{constant} (independent of $i$), which is too strong. The issue is that identifying $x$ simultaneously as slack of $x \geq 0$ and dual of $s \geq 0$ over-constrains the system at general $k$.

\subsection{Correct identification at $k = 1$}

At $k = 1$ (the initial barrier), the equations give $\sigma = 0$, so $d_{x,i} = -d_{s,i}$. More generally, the self-dual structure means the central path satisfies $x \circ s = \mu e$ and $\tau \kappa = \mu$, which in the geodesic parametrization becomes $d_x = -d_s$ and $d_\tau = -d_\kappa$ \textbf{at every point on the central path}.

To see this directly: on the central path,
\[
  x_i s_i = \mu, \quad \text{i.e.,} \quad \frac{(1 - d_{x,i})(1 - d_{s,i})}{k^2 w_{x,i} w_{s,i}} \cdot k^2 w_{x,i} w_{s,i} = \frac{(1-d_{x,i})(1+d_{x,i})}{1} = \mu
\]
Wait --- this is $x_i$ as slack times $\lambda_{x,i}$ as dual of the same constraint:
\[
  x_i \cdot \lambda_{x,i} = \frac{1 - d_{x,i}}{k w_{x,i}} \cdot k w_{x,i}(1 + d_{x,i}) = (1 - d_{x,i})(1+d_{x,i}) = 1 - d_{x,i}^2 = \mu.
\]
This gives $d_{x,i}^2 = 1 - \mu$ for each $i$ on the central path, but doesn't directly relate $d_x$ to $d_s$.

\subsection{Direct argument from the Newton equations}

The clearest argument: the Newton direction at $(W, k)$ solves a linear system. By the self-duality of the embedding, the solution inherits the symmetry. Specifically, if we apply the primal-dual swap $(x, s, \tau, \kappa) \to (s, x, \kappa, \tau)$ and $y \to -y$ to the Newton equations, and if $W$ satisfies $w_x = w_s^{-1}$ and $w_\tau = w_\kappa^{-1}$ (NT scaling for both pairs), then the transformed direction $(\tilde{d}_x, \tilde{d}_s) = (-d_s, -d_x)$ also solves the Newton equations. By uniqueness:
\[
  \boxed{d_x = -d_s, \qquad d_\tau = -d_\kappa.}
\]

\section{Implications for the KKT system}

\subsection{Reduced dimension}

With $d_s = -d_x$ and $d_\kappa = -d_\tau$, the direction has $n + 1$ independent cone components instead of $2n + 2$. The effective rank is
\[
  R_{\text{eff}} = n + 1.
\]

\subsection{Elimination of $s$ and $\kappa$}

Since $d_{s,i} = -d_{x,i}$, from \eqref{eq:slack}:
\[
  s_i = \frac{1 - d_{s,i}}{k w_{s,i}} = \frac{1 + d_{x,i}}{k w_{x,i}^{-1}} = k w_{x,i}(1 + d_{x,i}) = \lambda_{x,i}.
\]
So \textbf{$s = \lambda_x$}: the dual slack equals the dual variable of $x \geq 0$. Similarly $\kappa = \lambda_\tau$.

Substituting into the dual equality \eqref{eq:dual}:
\[
  -A^T y - \lambda_x + c\tau = r_d \theta.
\]
This is the standard dual feasibility condition with $s$ eliminated.

\subsection{Reduced KKT system}

After eliminating $s = \lambda_x$ and $\kappa = \lambda_\tau$, the independent variables are
\[
  (x, y, \tau, \theta) \in \R^{n + m + 2}
\]
with cone constraints $x \geq 0$, $\tau \geq 0$ (rank $n + 1$), and the equalities become:
\begin{align}
  Ax - b\tau &= r_p \theta, \\
  -A^T y - \lambda_x + c\tau &= r_d \theta, \\
  \ip{b}{y} - \ip{c}{x} - \lambda_\tau &= r_g \theta, \\
  \ip{r_p}{y} + \ip{r_d}{x} + r_g \tau &= -\alpha.
\end{align}
Here $\lambda_x = k w_x (1 + d_x)$ and $\lambda_\tau = k w_\tau (1 + d_\tau)$ are determined by the direction.

The Gram matrix reduces to $\diag(w_x^2, 0, w_\tau^2, 0)$ on the variables $(x, y, \tau, \theta)$, and the equality constraints involve $\lambda_x$ and $\lambda_\tau$ which couple back through $d_x$ and $d_\tau$.

\subsection{Practical impact}

\begin{itemize}
  \item The KKT system size drops from $(2n + m + 3) + (m + n + 2) = 3n + 2m + 5$ to $(n + m + 2) + (m + n + 2) = 2n + 2m + 4$. For the original LP with $m \ll n$, this nearly halves the system.
  \item The number of cone variables drops from $2n + 2$ to $n + 1$, halving the RowSpace.
  \item The $r$-updates (free back-solves) operate on $n + 1$ components instead of $2n + 2$.
  \item The complementarity gap is $\mu(R_{\text{eff}} - \|d\|^2) = \mu((n+1) - \|d_x\|^2 - d_\tau^2)$.
\end{itemize}

\end{document}
